    \def\stat{vikhrev}



\def\tit{О МЕХАНИЗМЕ РЕАЛИЗАЦИИ КОЭВОЛЮЦИОННОЙ МОДЕЛИ

ЖИЗНЕННОГО ЦИКЛА РАЗРАБОТКИ КОМПЬЮТЕРНЫХ ПРОГРАММ

ДЛЯ~ОБУЧЕНИЯ}



\def\titkol{О механизме реализации коэволюционной модели

жизненного цикла разработки} % компьютерных программ для~обучения}



\def\autkol{В.\,В.~Вихрев}



\def\aut{В.\,В.~Вихрев$^1$}



\titel{\tit}{\aut}{\autkol}{\titkol}



%{\renewcommand{\thefootnote}{\fnsymbol{footnote}}

%\footnotetext[1]

%{Работа выполнена при частичной финансовой поддержке  программы

%<<Интеллектуальные информационные технологии, системный анализ и

%автоматизация>> (проект 1.7).}}







\renewcommand{\thefootnote}{\arabic{footnote}}

\footnotetext[1]{Институт проблем информатики Российской академии наук, VVikhrev@ipiran.ru}





\label{st\stat}



                  \thispagestyle{headings}





     \Abst{На основе рассмотрения ситуации в ключевой отрасли

информатизации образования~--- отрасли разработки программного обеспечения (ПО)

для обучения~--- определяется круг проблем, стоящих перед отраслью,

и~предлагается подход к~решению нескольких, самых острых из них, путем

применения возможностей сервисов Интернета. Компьютерная программа~---

тот элемент, который превращает универсальное устройство

     ин\-фор\-ма\-ци\-он\-но-ком\-му\-ни\-ка\-ци\-он\-ных технологий (ИКТ)

в~средство обучения. Качество процесса обучения, зависящее от качества

программ, определяется, как следствие, состоянием дел в~отрасли разработки

ПО. Между тем эта отрасль исследована совершенно

недостаточно. Беглый анализ элементов дискурса участников процесса

информатизации образования показывает, что внешняя картина не

соответствует реальным отношениям внутри отрасли. Более глубокий анализ

приводит к~выводу, что начиная с~первых шагов информатизации отрасль

держалась в~значительной степени на инициативе одиночек и~небольших

команд энтузиастов. В~условиях хронического недофинансирования основным

ресурсом развития стала энергия энтузиазма. А~с~массовым включением

учителей в~процесс информатизации школы сложилась искаженная картина

о~месте учителя в~отрасли разработки ПО. На учителя

легла обязанность разработчика. Между тем его деятельность должна лишь

дополнять работу программиста. Правильно распределить обязанности

в~процессе разработки ПО для обучения поможет

система кооперации труда разработчиков и~учителей, основанная на модели

коэволюции, т.\,е.\ взаимосвязанного развития ПО

и~его применения в~конкретных педагогических технологиях. Для учителя эта

система предоставит консультативную информацию. Для разработчика система

станет средством дополнительного исследования предметной области

и~инструментом сопровождения, т.\,е.\ поддержкой тех этапов жизненного

цикла программы для обучения, которые в~наибольшей степени страдают от

недофинансирования. Система реализуется на основе сервисов Интернета.}



     \KW{информатизация образования; отрасль разработки программного

обеспечения; образовательный сегмент отрасли разработки программного

обеспечения; программное обеспечение для обучения; коэволюция;

коллективный разум}



\DOI{10.14357\!/\!08696527140411}







\vskip 12pt plus 9pt minus 6pt





      \thispagestyle{headings}



  \section{Введение}



    \vspace*{-2pt}



  Настоящая статья переводит в практическую плоскость проведенную ранее

аналитическую работу по исследованию процессов внедрения ИКТ в~школу

и~формирования на основе новых информационных технологий элементов

самоорганизации в~системе среднего образования~[1--8].



  Компьютерная программа~--- ключевое звено информатизации образования

в~той его части, что связана с~предметным преподаванием. Именно

специализированная компьютерная программа для обучения превращает

универсальные устройства ИКТ в~инструмент образования. За годы,

прошедшие с~начала информатизации образования в нашей стране (СССР,

затем Российская Федерация) сложилась отрасль разработки ПО

для образования. Она сложна и~достаточно слабо исследована.

Подобно тому, как финансирование образования не является приоритетной

статьей бюджетов разных уровней, эта отрасль также находится на втором

плане внутри финансового пространства системы образования. В~статье

обозначаются принципиальные особенности разработки ПО

 для обучения, формулируются главные проблемы отрасли

и~предлагается вариант решения части названных проблем в~сложившейся

инфраструктуре информационного образовательного пространства.



\vspace*{-4pt}



  \section{Отрасль разработки программного обеспечения

для~обучения и~ее~проблемы}



\vspace*{-4pt}



  \subsection{Что такое компьютерная программа и~программное

обеспечение для~обучения?}



  \vspace*{-2pt}



  Прежде всего ответим на вопрос: что означают термины \textit{программа}

и~\textit{программное обеспечение}? Ответ прост для любого, кто достаточно

близок к~предметной области <<информатика и~вычислительная техника>>.

Он может звучать так: программа~--- это набор команд, записанных на языке

того или иного уровня, совокупность которых (после перевода в~форму,

<<понятную>> компьютеру) реализует управление компьютером для

достижения ка\-кой-ли\-бо цели. Тогда программное обеспечение~--- это

совокупность программ, обеспечивающая возможность полезного

использования компьютера при решении ка\-ких-ли\-бо задач. Эти значения

возникли на заре вычислительной техники и~вряд ли изменятся в~обозримом

будущем.



  Компьютерная программа для обучения~--- это такая программа, которая

управляет компьютером для достижения ка\-кой-ли\-бо из целей обучения.

А~ПО для обучения~--- такая совокупность программ,

которая позволяет решать совокупность задач обучения. Под обучением здесь

и~далее понимается обучение предметное, т.\,е.\ изучение математики, физики,

истории, языков и~т.\,д.



  С определенного момента, однако, термина \textit{программа} оказывается

недостаточно для обозначения артефактов ПО.

Появляются термины \textit{пакет прикладных программ}, \textit{программный

комплекс}, \textit{программная сис\-те\-ма}. При взгляде на эти термины, даже не

пытаясь вникнуть в~их значения, далекий от программирования человек

способен сообразить, что сами по себе, набором входящих слов, они говорят об

определенных трансформациях в~об\-ласти разработки про\-грамм~--- об

увеличении в~размерах и~структурировании того, что называется программой,

например.



  О чем говорит появление терминов \textit{система управлением базой

данных} и~\textit{информационная система} или \textit{текстовый процессор}

и~\textit{издательская система}, заменяющих в дискурсе термин программа?

Первые два есть взгляд с~позиций программирования, они указывают, что

произошло качественное переосмысление некоего факта, известного

программистам и~ранее. А~именно: любая команда из набора, который выше

назван \textit{программой}, работает с~некоторыми \textit{данными}, несущими

\textit{информацию}. Но потребовалось проникновение вычислительной

техники в~новые области применения, в~сферы экономики и управления

социальными системами, в~частности, чтобы как откровение и~лозунг

прозвучало:

  про\-грам\-ма\;=\;струк\-ту\-ры данных\;+\;алго\-рит\-мы~\cite{9-vik}. Вторую

пару терминов, \textit{тексто\-вый процессор и издательская система}, можно

интерпретировать как взгляд на программу со стороны пользователя, когда

слово \textit{программа} заменяется в~обозначении ПО

названием главной цели или главной задачи его использования (языковый

прием метонимии).



  Трансформация терминосистемы во времени является индикатором

происходящих в~предметной области процессов даже для далекого от нее

наблюдателя. Взглянем с~этой точки зрения на изменения в~использовании

терминов терминосистемы ПО в~образовательном

дискурсе.



  Обратимся к материалам конференции <<Применение новых технологий

в~образовании>>, ежегодно проводимой в~г.~Троицке с~1990~г. Посмотрим,

какие термины использовали докладчики для обозначения того, что по своей

сути есть ПО для обучения, на трех конференциях,

отнесенных на десять лет одна от другой (выборка составлена по заголовкам

докладов):



\noindent

\begin{description}

\item[1994:] компьютерная программа, инструментальный пакет, исполнитель,

компьютер, персональный компьютер, про\-грам\-мно-ме\-то\-ди\-че\-ская

система, програм\-мно-ме\-то\-ди\-че\-ский комплекс, обучающая программа,

про\-грам\-мно-пе\-да\-го\-ги\-че\-ское средство (ППС), новая информационная

технология, компь\-ютер\-ная поддержка, инструментальное средство,

инструментальная система, обучающая среда, среда обучения, пакет программ,

компьютерное обеспечение, графический редактор, программная среда,

компьютерная лабораторная работа, обучающая контролирующая программа,

компьютерный учебный курс, игровой компьютерный комплекс;



\item[2004:] электронное учебное издание, компьютерная программа, программное

обеспечение, новая информационная технология, интерактивный курс,

образовательная среда на основе динамических слайд-лек\-ций,

мультимедийный урок-иг\-ра, компонент информационных технологий,

информационная сис\-те\-ма документооборота, информационная технология,

про\-грам\-мно-ме\-то\-ди\-че\-ский комплекс, информационная сис\-те\-ма,

компьютерная технология, интерактивная мультимедийная программа,

компьютерный тренажерный \mbox{комплекс}, сис\-те\-ма автоматизированного

проектирования, урок-те\-ле\-кон\-фе\-рен\-ция, презентационная технология,

моделирующая лабораторная работа, программный комплекс,

автоматизированный компьютерный комплекс, \mbox{e-learning} сис\-те\-ма, школьный

компьютерный тренажер, обучающая сис\-те\-ма, адаптивная обучающая сис\-те\-ма,

инструментальное интеллектуальное программное средство, электронная

презентация, виртуальная информационная среда;\\[-14pt]



\item[2014:] цифровой образовательный ресурс (ЦОР), электронный образовательный

ресурс (ЭОР), интернет-ресурс, ресурс он\-лайн-обуче\-ния, гипертекстовый

образовательный ресурс, сайт, электронное издание, электронный учебник,

электронное пособие, электронный учеб\-но-ме\-то\-ди\-че\-ский комплекс,

средство электронного обучения, образовательный продукт, программа,

компьютерная система, обуча\-ющая среда, обучающая видеоигра, виртуальная

игровая модель, авторское средство, мультимедийная лекция, цифровая

лаборатория, графический редактор, система динамической геометрии,

видеоквест, конструктор текстов, информационный ролик, компьютерная

модель, информационная система, средства ИКТ,

   ИКТ, информационная технология,  компьютерная технология.

\end{description}



  Сложность терминосистемы может свидетельствовать как о~сложности

предметной области, так и~о~сложных и~размытых отношениях между

акторами, действующими субъектами. С~позиций разработчика достаточно

легко прослеживается эволюция базового термина <<программа>> вслед за

эволюцией вычислительной техники:

\begin{description}

\item[1994:] программно-педагогический комплекс

  (про\-грам\-мно-пе\-да\-го\-ги\-че\-ское средство)~--- достаточно простые

персональные компьютеры, объединенные локальной сетью;



\item[2004:] электронное учебное издание (электронный учебник, мультимедийный

учебник)~--- реализация технологий мультимедиа и~ком\-пакт-диск как

средство доставки;



\item[2014:] электронный (цифровой) образовательный ресурс~--- единое

информационное пространство на основе Интернета.

\end{description}



  Совсем другое дело~--- пользователи ПО, в основном педагоги высшего

и~среднего образования. Основная масса приведенных терминов порождена

в~их дискурсе. Часть из них~--- своего рода адаптация терминологии

разработчиков, часть~--- продукт собственного понимания реалий предметного

поля. Бросается в~глаза обилие терминов, обозначающих, как следует из

докладов, ПО и~построенных по принципу синекдохи,

т.\,е.\ имя целого обозначает часть. Например, термин

  \textit{информационно-коммуникационная технология}  был

упо\-треб\-ля\-ем, как правило, именно в~случаях, когда имелось в виду ПО. Вот

фрагмент доклада учителя~\cite[с.~110]{10-vik}: <<По сравнению с~традиционной формой ведения

урока использование ИКТ высвобождает большое количество времени,

которое можно употребить для дополнительного объяснения материала. На

своих уроках в~целях развития творческих способностей я~использую

следующие виды ИКТ:

  \begin{itemize}

\item презентация~--- это наиболее распространенный вид представления

демонстрационных материалов;

\item обучающие игры и развивающие программы для младших

школьников;

\item дидактические материалы (представленные в электронном виде);

\item программы-тренажеры;

\item электронные учебники, энциклопедии>>.

\end{itemize}



  В известном смысле появление и укоренение термина электронный

(цифровой) образовательный ресурс можно рассматривать как бессилие

человеческого ума в~поисках ответа на вопрос: <<Что такое компьютерная

программа или ПО для обучения?>>



  \subsection{Краткая характеристика ситуации в~отрасли разработки

программного обеспечения для~обучения}



  Уже первичный анализ терминологии показывает, что отношения между

разработчиками ПО и~учителями, его потребителями,

отличаются от, так сказать, классических отношений

  про\-грам\-мист--поль\-зо\-ва\-тель. Они гораздо сложнее. Программа для

обучения для школы~--- это не просто товар, который можно выложить на

прилавок.



  Сложная оппозиция программист--учитель была замечена уже на заре

информатизации. Вот, например, четкая и~лаконичная формулировка данной

проблемы в~конце 1980-х~гг.~\cite[с.~6]{11-vik}: <<Программные средства учебного назначения

могут создаваться следующими способами:

  \begin{enumerate}[1.]

  \item Программное средство создается программистом~--- низкое качество с

точки зрения педагогики.

  \item Программное средство создается педагогом~--- низкое качество

с~точки зрения программиста.

  \item Программное средство создается при тесном контакте педагога

с~программистом (сюда же относится случай

  пе\-да\-го\-га-про\-грам\-миста)~--- качество высокое, но недостаточно

большое число таких коллективов.

  \item  Программисты создают некоторый пакет прикладных программ

(в~нашем понимании инструментальное средство), с помощью которого

  пе\-да\-гог-не\-про\-грам\-мист в~состоянии строить требующиеся

программные средства~--- возможность массовой разработки программных

средств учебного назначения, но значительные, в~смысле программирования,

сложности по разработке самих программных средств>>.

  \end{enumerate}



  Другими словами, решение оппозиции про\-грам\-мист--учи\-тель

предполагалось осуществить в рамках парадигмы <<автоформализации

знаний>> или <<персональных вычислений>> как этапа развития

информационных технологий. <<Основная задача персональных

  вычислений~--- формализации профессиональных знаний~--- выполняется,

как правило, полностью самостоятельно непрограммирующим профессионалом

или при минимальной технической поддержке программиста, который в~этом

случае имеет возможность включиться в~процесс формализации знаний только

на инструментальном уровне, оставляя наиболее трудную для его понимания

содержательную часть задачи специалисту в~данной предметной

  об\-ласти>>~\cite[с.~99]{12-vik}.



  Что же мы имеем в настоящее время? На поверхности картина выглядит

следующим образом. С~одной стороны, существует рынок ПО

 для обучения, товар на него поставляют несколько

софтверных компаний, к~которым в~последние годы присоединяются

издательства с~электронными приложениями к~тем учебникам, которые они

печатают. С~другой стороны, Интернет как информационное образовательное

пространство имеет в~бесплатном доступе хранилища

ЭОР (по сути~--- ПО для

обучения), где созданные профессионалами на государственные деньги

коллекции соседствуют с~коллекциями ресурсов социальных сообществ

учителей, материалами школьных и~учительских персональных сайтов.

Получается, что деятельность ко\-манд-раз\-ра\-бот\-чи\-ков гармонично

соединяется с~автоформализующими свои знания учителями. Однако

поверхностный взгляд не позволяет верно понять, ни кто по-на\-сто\-яще\-му

разрабатывает программы для обучения, ни в~каком состоянии находится

оппозиция раз\-ра\-бот\-чик--учи\-тель.



  Приглядимся к программистам разработчикам. Какие типовые команды

принимали участие в разработке в~процессе информатизации?



\smallskip



  \textbf{1988} (по данным~\cite{13-vik}):

  \begin{enumerate}[1.]

  \item  Специалисты специальных отделов академических институтов,

отраслевых институтов, на\-уч\-но-про\-из\-вод\-ст\-вен\-ных объединений,

вычислительных центров.

\item  Коллективы университетов, педагогических вузов, объединяющие

педагогов, инженеров и студентов.

  \item Коллективы из школ, объединяющие учителей и школьников.

  \item Независимые авторы, самостоятельно представляющие свою работу

(пользуясь формулой Ильфа и~Петрова~--- <<кус\-та\-ри-оди\-ночки

с~мотором>> (с~компьютером)).

  \end{enumerate}



\smallskip



\textbf{2008} (по данным об авторах Единой коллекции образовательных ресурсов):

  \begin{enumerate}[1.]

  \item  Специалисты отделов медиаиздательств и книжных издательств.

  \item Программисты коммерческих фирм.

  \item Коллективы кафедр вузов.



  \pagebreak



  \item Специалисты научно-исследовательских институтов.

  \item Независимые авторы.

  \end{enumerate}



  Анализ выявил дальнейшую градацию для кафедр вузов:

  \begin{enumerate}[(1)]

  \item <<россыпь талантов>>: институтские преподаватели с~научными

интересами в~той или иной области, приобщившиеся к~компьютерным

технологиям и~способные самостоятельно реализовать идеи по применению

компьютера в~обучении. Устойчивый интерес к~компьютерным технологиям

в~обучении реализуется в~форме своеобразного хобби. Основной результат

деятельности имеет форму фрагментарных разработок типа учебных модулей,

ЦОР;

  \item  <<россыпь талантов>> и~кафедра ИКТ: разработчики ЭОР по роду

профессиональной деятельности оказались далеки от компьютера в~смысле

личного владения навыками работы на нем. В~помощь привлекаются

сотрудники или лаборанты специализированных подразделений, например

кафедры ИКТ;

  \item  авторы учебника и разработчики набора ЦОР к~учебнику в~одном

вузе;

  \item  авторы учебника в вузе, ЦОР делают сторонние организации;

  \item  мультимедиалаборатории в вузе: специализированные подразделения

внут\-ри вуза, где преподаватели работают помимо основной деятельности на

кафедре;

  \item  мультимедиа-центры на стороне. Преподаватели реализуют свой

интерес к~компьютерам не в~рамках организационных структур внутри вуза,

а~на базе независимых структур, возможно, ими же созданных;

  \item  дистанционные профильные школы при вузах.

  \end{enumerate}



  К перечню 2008~г.\ необходимо добавить учителей и~школьников,

деятельность которых отчетливо проявляется в~материалах конференций по

вопросам информатизации образования. Таким образом, структура сообщества

разработчиков довольно сложна и, как показывает сравнение 1988 и 2008~гг.,

была таковой на всем протяжении информатизации.



  Ситуация в отрасли разработки ПО для обучения может быть в целом

охарактеризована следующими положениями:

  \begin{enumerate}[1.]

  \item  Систематическое, хроническое недофинансирование является даже не

проб\-ле\-мой, а~контекстом существования (жара в~пустыне или мороз в~тундре)

отрасли.

  \item  Рынок, на который выходят разработчики,~--- это не классический рынок

законченных товаров (то\-вар\;+\;ин\-струк\-ция на пятнадцати языках),

а~скорее рынок высокотехнологичного оборудования. Самоокупаемость здесь,

видимо, невозможна без мощных программ государственной поддержки.

  \end{enumerate}



  В таких суровых условиях выживают лишь энтузиасты и~увлеченные

бес\-среб\-реники, чья энергия и~является главным ресурсом развития отрасли

ПО для обучения. Приходят годы тучные, появляется

государственный заказ, и~вокруг них группируются команды из энтузиаста

(или группы энтузиастов) и~временных помощников. Приходят годы тощие,

и~они снова в~одиночестве продолжают свою творческую работу. Часто они

действуют сами по себе, но, как правило, их привлекают к~работе фирмы

и~издательства, выпуская в свет их труд под своим торговым знаком.



  Чтобы энтузиазм не выгорел дотла, он должен быть поддержан. А~для этого

необходимо понять, какого рода проблемы нужно решить в~отрасли разработки

ПО для обучения.



  \subsection{Каноническая схема жизненного цикла программного

обеспечения}



  Чтобы сформулировать представление о проблемах, необходимо выбрать

точку отсчета. В~качестве такой точки отсчета или нормы, которой можно

поверить ситуацию в~отечественной отрасли разработки

ПО для обучения, возьмем концепцию жизненного цикла программы.



  Считается, что до начала 1980-х~гг.\ программирование было скорее

ре\-мес\-лом и искусством, чем промышленной профессией~[14]. Однако

громадный рост потребности в программах и~повышение цены ошибки

программиста по мере нарастания интенсивности процесса информатизации

сначала промышленности, экономики, а~затем и~всего общества, потребовал

пересмотра взглядов на организацию труда программистов. Появились

и~начали регулярно обновляться и~переиздаваться труды по научной

организации процесса разработки ПО. Оформилось представление

о~жизненном цикле программы, который включает следующие этапы:

  \begin{enumerate}[(1)]

  \item  постановка задачи, формулировка требований;

  \item  исследование проблемной области, разработка проекта;

  \item  программирование, отладка;

  \item  тестирование, доработка по результатам тестирования;

  \item  сопровождение, доработка в процессе эксплуатации;

  \item  формирование требований на переработку.

  \end{enumerate}



  Хотя внимание большинства авторов монографий и~учебников~[15, 16],

посвященных разработке ПО, сосредоточено на первых этапах жизненного

цикла, существует также обширная литература, ориентированная на

особенности этапа сопровождения~[17]. Этот этап признается столь важным,

что предлагается выделять в~команде специалиста, ответственного

исключительно за сопровождение. В~процессе сопровождения в~канонической

схеме происходит накопление требований к~обновлению программы.



  Стремительное развитие вычислительной техники в~эпоху персональных

компьютеров трансформировало исходную схему жизненного цикла. Теория

эволюции ПО~[18, 19], прежде применявшаяся

к~большим программным сис\-те\-мам, теперь может быть распространена

практически на все ПО. Эволюция, т.\,е.\ выпуск обновленных версий старой

программы, происходит по нескольким ключевым причинам. Обновившаяся

техническая база стимулирует непрерывное углубление представления

о~проблемной области, для которой предназначена программа, что приводит

к~реализации новых функций в~программе. При неизменных функциях может

совершенствоваться интерфейс взаимодействия с~пользователем. Наконец, при

неизменных функциях и~неменяющемся интерфейсе могут оптимизироваться

ресурсы программы. Все это означает, что в~работу по сопровождению должен

включаться и~пользователь программы, которому необходимо отслеживать

происходящие изменения и~проецировать их на свою деятельность. Другими

словами, по мере становления информационного общества и~роста

информационной и~технокультуры элементом этой культуры становится

интерес к~новостям в~той области ПО, на которой

строится информационная работа пользователя.



  Необходимо отметить, что, поскольку каноническая схема в~целом

ориентирована на программу как товарный продукт, в~ней не предусмотрено

процедур систематического взаимодействия пользователей и~программистов.

Такие взаимодействия хотя и~подразумеваются, однако имеют характер

отдельных акций.



  \section{Формулировка проблем}



  Как было отмечено, отрасль разработки ПО для

обучения функционирует в~условиях хронического недофинансирования. При

этом обучение является одним из сложнейших видов деятельности.

Педагогическая наука пока не исследовала его до такой степени, чтобы

относительно просто различать содержательную и~дидактическую сторону

этого процесса. Всякий раз соединение учебного материала со способом его

подачи является искусством. В~этом ключевая проблема сложности

предметной области. Кроме того, разнообразие условий обучения и~различие

качества участников этого процесса делает неразрешимой задачу однозначного

методического описания. Как следствие~--- спектр проблем разработчика

и~поль\-зо\-ва\-теля-учи\-теля.



  Проблемы разработчика:

  \begin{enumerate}[1.]

\item Точность постановки задачи.

\item Разработка грамотного методического обеспечения.

\item Привязка ситуации к разнообразным ситуациям и~условиям

применения.

\item Сопровождение, наличие сил и возможностей для грамотного

решения за\-дачи.

  \end{enumerate}



  Главная проблема разработчика ПО для обучения~---

он не обладает ресурсами для полноценной реализации канонической схемы

жизненного цикла.



  Проблемы учителя:

  \begin{enumerate}[1.]

\item Овладение возможностями ПО при недостаточной помощи

разработчика (у~которого на эту помощь нет ресурсов).

\item Адаптация возможностей ПО к конкретной ситуации в~силу

недостаточной собственной квалификации или особой сложности

ситуации.

\item Разработка педагогической технологии на основании использования

возможностей имеющейся программы.

  \end{enumerate}



  Главная проблема учителя~--- он не обладает ресурсами для детального

и~достаточно самостоятельного углубления в~программу и~ее возможности.



  Заметим, что при нехватке материальных ресурсов ситуация отчасти

ис\-прав\-ля\-ет\-ся, а~проблемы решаются за счет главного человеческого

  ресурса~--- времени жизни.



  И, наконец, в условиях ограниченности материальных ресурсов необходимо

компенсировать работу энтузиаста хотя бы тем, что он получает

удовлетворение от вида результатов своего труда~--- проблема <<обратной

связи>>, важная не только для разработчика, но и~для пользователя.



  \section{Коэволюционная схема жизненного цикла программного

обеспечения}



  \subsection{Разрешение оппозиции программист--учитель}



  Предположение о том, что учитель должен (или может) выступать в роли

про\-грам\-миста-раз\-ра\-бот\-чи\-ка, как кажется, подтверждается жизненной

практикой. Количество ЭОР, которые

учителя выложили в~Интернет, едва ли уступает, а~возможно, и~превосходит

количество ресурсов, созданных профессиональными разработчиками. Однако

этот факт, на взгляд автора, лишь затемняет истинное положение вещей.

Несомненно, существуют учителя, и~их немало, которые на вполне

профессиональном уровне освоили ремесло разработчика. Они могут создавать

вполне конкурентоспособные программы. Это не меняет сути дела, которая

состоит в~том, что учитель, включаясь в~процесс информатизации, входит

туда со своей особой ролью.



  О правомерности использования термина \textit{педагогическая технология}

дискуссия в~педагогическом сообществе еще продолжается. У~него много

противников. В~данном случае важно, что технология~--- универсальный

способ описания сложного процесса по шагам. Тогда педагогическая

технология~--- описание педагогического процесса или процесса обучения.

Задача учителя состоит в~том, чтобы спроектировать этот процесс

с~включением в~него ИКТ. Фактически педагогический процесс соединяется

с~ин\-фор\-ма\-ци\-он\-но-ком\-му\-ни\-ка\-ци\-он\-ным процессом,

поддержанным техническими, программными и информационными ресурсами.



  В идеале учитель привлекает некоторое ПО для обучения в~качестве

информационной системы, основы процесса. Но даже теоретически не

существует такого ПО, которое идеально вписалось бы в~<<ландшафт>> урока.

Необходима привязка к~местности. Такая привязка может быть названа

адаптацией. Существует три типа адаптации:

  \begin{enumerate}[(1)]

  \item параметрическая~--- путем настройки параметров;

  \item модульная~--- путем отбора используемых блоков;

  \item конструктивная~--- путем разработки недостающих блоков.

  \end{enumerate}







  \begin{figure}[b] %fig1

  \vspace*{1pt}

 \begin{center}

 \mbox{%

 \epsfxsize=86.659mm

 \epsfbox{vih-1.eps}

 }

 \end{center}

 \vspace*{-9pt}

  \Caption{Отражение программы для обучения (информационной системы)

  в~ин\-фор\-ма\-ци\-он\-но-обра\-зо\-ва\-тель\-ное пространство}

  \end{figure}



  При огромном дефиците качественного ПО для

обучения учитель невольно оказывается в~ситуации, когда он сам вынужден

конструировать, разрабатывать программы. Большинство созданных учителями

ЭОР вызвано необходимостью

конструктивной адаптации существующего ПО.



  Другая часть ЭОР, внесенных

  в~ин\-фор\-ма\-ци\-он\-но-обра\-зо\-ва\-тель\-ное пространство учителями,

принадлежит к~типу <<электронная версия методики проведения урока>>, т.\,е.\

технологическая карта педагогической технологии. Именно эти ЭОР

представляют особую ценность. Любое ПО

  в~ин\-фор\-ма\-ци\-он\-но-обра\-зо\-ва\-тель\-ном пространстве можно

представить окруженным такими педагогическими электронными модулями

(рис.~1).



  \subsection{Анализ ситуации разработки программного обеспечения

сквозь~призму~концепции коллективного разума}



  Описанная ситуация идеально ложится на концепцию коллективного

интеллекта~[20, 21].



  Разработчик ПО для обучения и учителя, применяющие в~своей практике его

разработку, образуют социальную группу с~четким разделением социальных

ролей и~хорошо понятными ролевыми функциями. Хотя каждый из участников

действует автономно, у~них есть интерес объединиться для интеллектуального

исследования ситуации.



  При адекватном ситуации поведении участников их коммуникация позволит

выявлять проблемы и~находить пути их решения, т.\,е.\ реализовать основную

функцию интеллекта. Так возникает коллективный разум.



\begin{figure}[b] %fig2

  \vspace*{-6pt}

 \begin{center}

 \mbox{%

 \epsfxsize=78.318mm

 \epsfbox{vih-2.eps}

 }

 \end{center}

 \vspace*{-9pt}

\Caption{Общая структура информационной системы поддержки коэволюционной модели

разработки ПО (ИСП КМР ПО)}

\end{figure}



\vspace*{-4pt}





  \subsection{Предпосылки для развития коммуникации}



  \vspace*{-2pt}



  Все стороны, участвующие в коммуникации, заинтересованы в развитии на

основе этой коммуникации. Возникают предпосылки для совместной

эволюции~--- коэволюции.



  Разработчик, по существу, реализует функцию этапа сопровождения,

одновременно исследуя проблемную область (уточнение постановки задачи на

ПО) и~отслеживая потребности в обновлении версии (эволюция ПО).



  Для учителя\,--\,пользователя ПО существует объективная потребность

в~обмене методическим опытом. В~рамках школы учителя обычно образуют

методические объединения, которые часто выходят за рамки школы.

Материалы конференций показывают, что уже на ранних шагах

информатизации она рассматривалась как средство накопления опыта (базы

методических данных), налаживания контактов между заинтересованными

в~методическом общении. Развитие Интернета и~формирование единого

информационного образовательного пространства сопровождалось

формированием в~нем методического сегмента.





  При этом основной поток коммуникаций осуществляется через

информационную систему поддержки коэволюционной модели разработки ПО

(связи, обозначенные буквой~А на рис.~2), что не исключает возможности

контактов по другим каналам (связи~Б).





  \subsection{Реализация модели коэволюционного развития

программного обеспечения для обучения}



  Основой системы является сайт разработчика ПО.

Его функциональная архитектура представлена на рис.~3.



\begin{figure}[h] %fig3

  \vspace*{1pt}

 \begin{center}

 \mbox{%

 \epsfxsize=127.825mm

 \epsfbox{vih-3.eps}

 }

 \end{center}

 \vspace*{-9pt}

\Caption{Функциональная архитектура ИСП КМР ПО}

\end{figure}





  Обязанности модератора ложатся на разработчика как на энтузиаста

продвижения своего ПО. Стандартный режим работы:

  \begin{itemize}

  \item привлечение новых участников;\\[-14pt]

  \item обмен мнениями;\\[-14pt]

  \item информационная рассылка;\\[-14pt]

  \item выкладывание методик;\\[-14pt]

  \item выкладывание новых версий ПО.

  \end{itemize}



  \vspace*{-6pt}



  \section{Заключение}



    \vspace*{-2pt}



  Информатизация образования~--- непрерывно разворачивающийся процесс.

Для сверхсложной социотехнической системы, каковой становится система

образования по мере информатизации, управляющее регулирование сверху

должно сочетаться с~саморегулированием и~саморганизацией.

Самоорганизация держится на деятельности малой части социального

сообщества, на людях, склонных к~повышенной активности, а~также к~особой

форме проявления активности~--- к~инновациям.

Должны создаваться

механизмы, ячейки коллективного разума, способные поддерживать развитие

системы на ресурсе энергии человеческой инициативы. Задача системы

регулирования~--- подпитывать точки роста другими видами ресурсов.

Предложенный в~статье механизм самоорганизации в~рамках отрасли

разработки ПО для обучения может быть назван такой

точкой.



  \vspace*{-4pt}



{\small\frenchspacing

{ %\baselineskip=10pt

\addcontentsline{toc}{section}{Литература}

\begin{thebibliography}{99}



  \vspace*{-2pt}



\bibitem{1-vik}

\Au{Вихрев В.\,В.} Смена парадигмы: от электронного учебника к

ЦОРу~/\!\!/ Новые образовательные технологии в~вузе: Сб. мат-лов

VII~Междунар. на\-уч.-ме\-то\-ди\-ч. конф.~--- Екатеринбург:

УГТУ-УПИ им.\ Б.\,Н.~Ельцина, 2010.  Ч.~2. С.~20--24.

\bibitem{2-vik}

\Au{Вихрев В.\,В., Шпакова Т.\,Ю.} Компьютерное творчество учителей

как ресурс информатизации образования~/\!\!/ Новые образовательные

технологии в~вузе: Сб. мат-лов VII Междунар. на\-уч.-ме\-то\-дич.

конф.~--- Екатеринбург: УГТУ-УПИ им. Б.\,Н.~Ельцина, 2010.  Ч.~2.

С.~24--28.

\bibitem{3-vik}

\Au{Вихрев В.\,В.} О методической поддержке учителя, ключевой фигуры

информатизации школы~/\!\!/ Информационные технологии для Новой

школы: Мат-лы конф.~--- СПб: Региональный центр оценки качества

образования и~информационных технологий, 2010. С.~92--95.

\bibitem{4-vik}

\Au{Вихрев В.\,В.} Концептуальная модель представления системы

образования как аутопоэтической системы~/\!\!/ Всеросс. конф.

<<Информационные технологии в~образовании XXI~века>>. Всеросс.

молодежная науч. конф. <<Информационные технологии

в~образовательном процессе исследовательского университета>>: Сб.

науч. тр.~--- М.: НИЯУ МИФИ, 2011. С.~381--385.

\bibitem{5-vik}

\Au{Вихрев В.\,В.} О некоторых аспектах влияния технологических

изменений на программное обеспечение для образования~/\!\!/

Применение ЭОР в~образовательном

 процессе (ИTO-ЭОР-2011): Тезисы

докл. Всеросс. конф.~--- Обнинск: ИАТЭ НИЯУ МИФИ, 2011. С.~50--51.



%\pagebreak



\bibitem{6-vik}

\Au{Вихрев В.\,В.} О концептуальной модели информатизации системы

образования (в~порядке постановки задачи)~/\!\!/ Современные

информационные технологии и~ИТ-обра\-зо\-ва\-ние: Сб. избранных докл.

VI Междунар. на\-уч.-прак\-тич. конф.~--- М.: ИНТУИТ.РУ, 2011. С.~111--120.

\bibitem{7-vik}

\Au{Вихрев В.\,В.} Электронный образовательный ресурс: через

феноменологию к типологии~/\!\!/ Ученые записки ИСГЗ, 2013. №\,1(11).

С.~72--79.

\bibitem{8-vik}

\Au{Вихрев В.\,В.} Инь и ян процесса информатизации (к феноменологии

термина <<электронный образовательный ресурс>>)~/\!\!/ Применение

инновационных технологий в~образовании (ИТО-Троицк-2013): Сб. докл.

XXIV Междунар. конф.~---  Троицк--Москва: Байтик, 2013. С.~163--166.

\bibitem{9-vik}

\Au{Вирт Н.} Структуры данных\;+\;алгоритмы\;=\;про\-грам\-ма~/ Пер.

с~англ.~--- М.: Мир, 1985. 406~с.

(\Au{Wirth N.} Algorithms\;+\;data structures\;=\;programs.~---

Englewood Cliffs, NJ: Prentice-Hall, Inc., 1976. 366~p.)

\bibitem{10-vik}

\Au{Бабичек И.\,А.} Развитие способностей младших школьников во

внеурочной деятельности через использование

ин\-фор\-ма\-ци\-он\-но-ком\-му\-ни\-ка\-ци\-он\-ных технологий~/\!\!/

Применение новых технологий в~образовании: Мат-лы XXV

Междунар. конф.~--- Троицк--Москва: Байтик, 2014. 640~с.

\bibitem{11-vik}

\Au{Федосеев А.\,А., Христочевский С.\,А.} Об инструментальных

средствах для разработки ППС~/\!\!/ Разработка и применение

программных средств ПЭВМ в учебном процессе: Мат-лы IV Всесоюзн.

семинара.~--- М.: ИПИАН, 1988. С.~6--8.

\bibitem{12-vik}

\Au{Громов Г.\,Р.} Национальные информационные ресурсы: проблемы

промышленной эксплуатации.~--- М.: Наука, 1985. 240~с.

\bibitem{13-vik}

Компьютеризация образования (средства, методы, передовой опыт):

Тематическая выставка.~--- М.: ВДНХ, 1988. 109~с.

\bibitem{14-vik}

\Au{Боэм Б.\,У.} Инженерное проектирование программного

обеспечения~/ Пер. с~англ.~--- М.: Радио и связь, 1985. 512~с.

(\Au{Boemh B.\,U.} {Software engineering economics}.~---

      Englewood Cliffs, NJ: Prentice Hall, Inc., 1981. 767~p.)

\bibitem{15-vik}

\Au{Брауде Э.\,Дж.} Технология разработки программного обеспечения~/

Пер. с~англ.~--- СПб.: Питер, 2004. 655~с. (\Au{Braude~E.\,J.} Software

engineering: An object-oriented perspective.~--- N.Y.: John Wiley\,\&\,Sons,

2001. 560~p.).

\bibitem{16-vik}

\Au{Гагарина Л.\,Г., Кокорева Е.\,В., Виснадул~Б.\,Д.} Технология

разработки про\-грам\-мно\-го обеспечения.~--- М.: Форум\,--\,ИНФРА-М,

2008. 400~с.

\bibitem{17-vik}

\Au{Гласс Р., Нуазо Р.} Сопровождение программного обеспечения~/ Пер.

с~англ.~--- М.: Мир, 1983. 156~с. (\Au{Glass~R.\,L., Noiseux~R.\,A.}

Software maintenance guidebook.~--- Englewood Cliffs, NJ: Prentice-Hall, Inc., 1981.

204~p.).

\bibitem{18-vik}

\Au{Hongji Y., Ward M.} Successful evolution of software systems.~---

Norwood: Artech House, 2003. 282~p.

\bibitem{19-vik}

Software evolution~/ Eds.\ T.~Mens, S.~Demeyer.~--- Berlin, Heidelberg: Springer-Verlag, 2008. 347~p.



\bibitem{21-vik} %20

\Au{Шуровьески Дж.} Мудрость толпы. Почему вместе мы умнее, чем

поодиночке, и~как коллективный разум влияет на бизнес, экономику,

общество и~государство~/ Пер. с~англ.~--- М.: Вильямс, 2007. 304~с.

(\Au{Surowiecki J.} The wisdom of crowds. Phillips\,\&\,Nelson Media, 2004. 336~p.).



\bibitem{20-vik} %21

Collective intelligence: Creating a prosperous world at peace~/ Ed.\

M.~Tovey.~--- Oakton: Earth Intelligence Network, 2008. 612~p.

\end{thebibliography}

} }





\vspace*{-6pt}



\hfill{\small\textit{Поступила в редакцию 18.09.14}}



%\newpage





\vspace*{7pt}



\hrule



\vspace*{2pt}



\hrule



%\vspace*{9pt}





\def\tit{ON THE MECHANISM OF IMPLEMENTATION OF~COEVOLUTIONARY MODEL

     OF~THE~DEVELOPMENT LIFE~CYCLE OF~COMPUTER PROGRAMS FOR~LEARNING}







\def\titkol{On the mechanism of implementation of~coevolutionary model

     of~the~development life cycle} % of~computer programs for~learning}



\def\aut{V.\,V.~Vikhrev}

\def\autkol{V.\,V.~Vikhrev}





\titel{\tit}{\aut}{\autkol}{\titkol}



\vspace*{-12pt}



\noindent

Institute of Informatics Problems, Russian Academy of Sciences,

44-2 Vavilov Str., Moscow 119333, Russian

Federation







\def\leftfootline{\small{\textbf{\thepage}

\hfill Sistemy i Sredstva Informatiki~---

Systems and Means of Informatics 2014 vol~24 no\,4}

}%

 \def\rightfootline{\small{Sistemy i Sredstva Informatiki~---

 Systems and Means of Informatics   2014   vol~24   no\,4

\hfill \textbf{\thepage}}}





     \Abste{The article considers the situation in the key

     sector of education informatization (i.\,e., the field of education

     software development), determines

     the range of problems faced by the

     main sector, and suggests an appropriate approach to solving the most serious

     problem by using Internet services.

A~computer program is an element that turns a universal device of information

and\linebreak\vspace*{-12pt}}



     \Abstend{communication technologies into an education tool. The quality of training

process depends on the quality of programs and, consequently, is defined by

the state of affairs in the field of software development. Meanwhile,

this sector of education informatization is studied poorly.

A~quick analysis of the discourse elements of participants in the

education informatization process demonstrates that the external picture

does not correspond to the real relations within the sector. An astute

analysis leads to the conclusion that from the very beginning of informatization,

the industry has kept largely on the initiative of individuals and small teams

of enthusiasts.

Under the conditions of chronic underfunding, the energy of enthusiasm has

become a major resource for development. The active inclusion of teachers

into the process of school informatization has created a distorted picture

of teacher's place in the industry of software development. The function

of development engineer has been placed on teachers. But their activities

should only supplement the work of programmers.

Therefore, the cooperative system of development for engineers' and teachers'

labour based on the coevolutionary model (i.\,e., the interrelated software

development and its application in the specific pedagogical technologies)

can help allocate duties right in the process of education software

development. For teachers, this system can provide advisory data. For

development engineers, the system can become a~means of a further research

in the subject area and can become a~support tool of the stages of education

software development life cycle that primarily suffer from underfunding.

The system is implemented on the basis of Internet services.}



     \KWE{informatization of education; software development; educational

software; coevolution; collective intelligence}







\DOI{10.14357\!/\!08696527140411}



%\vspace*{-24pt}



%\Ack

%\noindent





\vspace*{-12pt}







\renewcommand{\bibname}{\protect\rmfamily References}

%\renewcommand{\bibname}{\large\protect\rm References}



{\small\frenchspacing

{%\baselineskip=10.8pt

\addcontentsline{toc}{section}{References}

\begin{thebibliography}{99}

\bibitem{1-vik-1}

\Aue{Vikhrev, V.\,V.} 2010. Smena paradigmy: Ot elektronnogo uchebnika

k~TsORu [Paradigm shift: From an electronic textbook to a digital education

resource]. \textit{Novye Obrazovatel'nye Tekhnologii v~Vuze: Sb. mat-lov

VII~Mezhdunar. Nauch.-Metodich. Konf.} [New Educational Technologies in

Higher Education: 7th  Scientific-Methodical Conference (International)

Proceedings]. Ekaterinburg. 2:20--24.

\bibitem{2-vik-1}

\Aue{Vikhrev, V.\,V., and T.\,Yu.~Shpakova}. 2010. Komp'yuternoe tvorchestvo

uchiteley kak resurs informatizatsii obrazovaniya [Computer creativity of teachers as

a resource of informatization of education]. \textit{Novye Obrazovatel'nye

Tekhnologii v~Vuze: Sb. mat-lov VII Mezhdunar. Nauch.-Metodich. Konf.}

[New Educational Technologies in Higher Education: 7th Scientific-Methodical

Conference (International) Proceedings]. Ekaterinburg. 2:24--28.

     \bibitem{3-vik-1}

     \Aue{Vikhrev, V.\,V.} 2010. O~metodicheskoy podderzhke uchitelya,

klyuchevoy figury informatizatsii shkoly [On the methodological support of teachers,

key of school informatization]. \textit{Informatsionnye Tekhnologii dlya Novoy

Shkoly: Mat-ly Konf}. [Information Technology for the New School: Conference

Proceedings]. St.\  Petersburg. 92--95.

\bibitem{4-vik-1}

\Aue{Vikhrev, V.\,V.} 2011. Kontseptual'naya model' predstavleniya sistemy

obrazovaniya kak autopoeticheskoy sistemy [Conceptual model of representation of

the education system as an autopoietic system]. \textit{Vseross. Konf.

``Informatsionnye Tekhnologii v~Obrazovanii XXI~Veka.'' Vseross.

Molodezhnaya Nauch. Konf. ``Informatsionnye Tekhnologii v Obrazovatel'nom

Protsesse Issledovatel'skogo Universiteta:'' Sb. nauch. tr}. [All-Russian

Conference ``Information Technologies in Education of the XXI Century.''

All-Russian Youth Scientific Conference ``Information Technologies in the

Educational Process Research University: Collection of scientific works].  Moscow.

381--385.

     \bibitem{5-vik-1}

     \Aue{Vikhrev, V.\,V.} 2011. O nekotorykh aspektakh vliyaniya

tekhnologicheskikh izmeneniy na programmnoe obespechenie dlya obrazovaniya [On

some aspects of the impact of technological changes on software for education].

\textit{Primenenie EOR v~Obrazovatel'nom Protsesse (ITO-EOR-2011): Tezisy

dokl. Vseross. Konf.} [The Use of e-Learning Resources in the Educational

Process (ITO-RAR-2011): Abstracts of the All-Russian Conference]. Obninsk. 50--51.

     \bibitem{6-vik-1}

     \Aue{Vikhrev, V.\,V.} 2011. O kontseptual'noy modeli informatizatsii sistemy

obrazovaniya (v~poryadke postanovki zadachi) [Conceptual model of informatization

of the education system (in the order of statement of the problem)].

\textit{Sovremennye Informatsionnye Tekhnologii i~IT-Obrazovanie: Sb. izbrannykh

dokl. VI~Mezhdunar. Na\-uch.-Prak\-tich. Konf.}

[Modern Information Technology and

IT-Education: A~Collection of Selected Papers of the 6th

     Scientific-Practical Conference (International). Moscow. 111--120.

     \bibitem{7-vik-1}

     \Aue{Vikhrev, V.\,V.} 2013. Elektronnyy obrazovatel'nyy resurs: Cherez

fenomenologiyu k tipologii [Electronic educational resource: Through

phenomenology to typology]. \textit{Uchenye Zapiski ISGZ} [Kazan ISGS ``Scientific

Notes''].  1(11):72--79.

\bibitem{8-vik-1}

\Aue{Vikhrev, V.\,V.} 2013.  In' i yan protsessa informatizatsii (k~fenomenologii

termina ``elektronnyy obrazovatel'nyy resurs'') [Yin and Yang of the process of

informatization (on phenomenology of the term ``electronic educational resource'')].

\textit{Primenenie Innovatsionnykh Tekhnologiy v~Obrazovanii (ITO-Troitsk-2013):

Sb. dokl. XXIV~Mezhdunar. Konf.} [New Computer Technology in Education:

24th  Conference (International) Proceedings]. Moscow--Troitsk:

Bytik. 163--166.

     \bibitem{9-vik-1}

     \Aue{Wirth, N.} 1976. \textit{Algorithms\;+\;data structures\;=~programs}.

Englewood Cliffs, NJ: Prentice-Hall, Inc. 366~p.

     \bibitem{10-vik-1}

     \Aue{Babichek, I.\,A.} 2014. Razvitie sposobnostey mladshikh shkol'nikov vo

vneurochnoy deyatel'nosti cherez ispol'zovanie in\-for\-ma\-tsi\-on\-no-kom\-mu\-ni\-ka\-tsi\-on\-nykh

tekhnologiy [The development of the abilities of younger students in extracurricular

activities through the use of information and communication technologies].

\textit{Primenenie Novykh Tekhnologiy v~Obrazovanii: Mat-ly XXV Mezhdunar.

Konf.} [New Computer Technology in Education: 25th Conference (International)

Proceedings]. Moscow--Troitsk: Bytik. 109--111.

\bibitem{11-vik-1}

\Aue{Fedoseev, A.\,A., and S.\,A.~Khristochevskiy}. 1988. Ob instrumental'nykh

sredstvakh dlya razrabotki PPS [On tools for developing software and educational

tools]. \textit{Razrabotka i~Primenenie Programmnykh Sredstv PEVM v~Uchebnom

Protsesse: Mat-ly IV~Vsesoyuzn. Seminara.} [Development and Application of

Software PC Tools in the Educational Process: 4th All-Union  Seminar

Proceedings]. Moscow. 6--8.

     \bibitem{12-vik-1}

     \Aue{Gromov, G.\,R.} 1985. \textit{Natsional'nye informatsionnye resursy:

Problemy promyshlennoy ekspluatatsii} [National information resources: Problems

of industrial exploitation]. Moscow: Nauka. 240~p.

     \bibitem{13-vik-1}

\textit{Komp'yuterizatsiya obrazovaniya (sredstva, metody, peredovoy

opyt): Tematicheskaya vystavka} [Computerization of education (means, methods,

best practices): A~thematic exhibition]. 1988. Moscow: VDNKh. 109~p.

     \bibitem{14-vik-1}

     \Aue{Boemh, B.\,U.} 1981. \textit{Software engineering economics}.

      Englewood Cliffs, NJ: Prentice Hall, Inc. 767~p.

     \bibitem{15-vik-1}

     \Aue{Braude, E.\,J.} 2001. \textit{Software engineering: An object-oriented

perspective}.  N.Y.: John Wiley\,\&\,Sons. 560~p.

     \bibitem{16-vik-1}

     \Aue{Gagarina, L.\,G., E.\,V.~Kokoreva, and B.\,D.~Visnadul}. 2008.

\textit{Tekhnologiya razrabotki programmnogo obespecheniya} [The technology of

software development]. Moscow: Forum\,--\,INFRA-M. 400~p.

     \bibitem{17-vik-1}

     \Aue{Glass, R.\,L., and R.\,A.~Noiseux}. 1981. \textit{Software maintenance

guidebook}.  Englewood Cliffs, NJ: Prentice Hall, Inc. 204~p.

     \bibitem{18-vik-1}

     \Aue{Hongji, Y., and M.~Ward}. 2003. \textit{Successful evolution of

software systems.}  Norwood: Artech House. 282~p.

     \bibitem{19-vik-1}

     Mens, T., and S.~Demeyer, eds. 2008. \textit{Software evolution}.  Berlin,

Heidelberg: Springer-Verlag. 347~p.



     \bibitem{21-vik-1}

     \Aue{Surowiecki, J.} 2004. \textit{The wisdom of crowds}.

Phillips\,\&\,Nelson Media. 336~p.



     \bibitem{20-vik-1}

     Tovey, M., ed. 2008. \textit{Collective intelligence: Creating a~prosperous

world at peace}. Oakton: Earth Intelligence Network. 612~p.

\end{thebibliography}

} }





\vspace*{-6pt}



\hfill{\small\textit{Received September 18, 2014}}



     \Contrl



     \noindent

     \textbf{Vikhrev Vladimir V.} (b.\ 1957)~--- senior scientist, Institute of

Informatics Problems, Russian Academy of Sciences, 44-2 Vavilov Str., Moscow

119333, Russian Federation; VVikhrev@ipiran.ru





 \label{end\stat}



\renewcommand{\bibname}{\protect\rm Литература}